% Partie : applied-modal-logics
% Chapitre : epistemic-logic
% Section : properties-accessibility

\documentclass[../../../include/open-logic-section]{subfiles}

\begin{document}

\olfileid[fr]{aml}{el}{acc}

\olsection{Relations d'accessibilité et principes épistémiques}

À partir de ce que nous savons déjà de la correspondance entre
formules et cadres en logique modale normale, nous pouvons examiner
les !!{formula}s caractéristiques sous une interprétation épistémique.
Nous avons indiqué que les logiques épistémiques sont souvent
interprétées dans S5. Voyons donc comment divers principes épistémiques
sont représentés et quelles conditions sur les cadres leur correspondent.

Rappelons qu'en logique modale normale, différentes !!{formula}s
modales caractérisent différentes propriétés des relations
d'accessibilité. Le tableau suivant présente quelques principes
épistémiques et les conditions correspondantes.

\begin{table}[t]
    \begin{tabular}{| p{\dimexpr.5\linewidth-2\tabcolsep-2\arrayrulewidth-.5\doublerulesep\relax} || p{\dimexpr.5\linewidth-2\tabcolsep-2\arrayrulewidth-.5\doublerulesep\relax} |}
      \hline
      {\emph{Si $R$ est \dots}} & {\emph{alors \dots est vraie en tout monde de~$\mModel{M}$ :}} \\
      \hline \hline

      quelconque & $\Knows (p \lif q) \lif (\Knows p \lif \Knows q)$
      \hfill \newline{(Clôture)} \\
      \hline
      \emph{réflexive} : $\forall w Rww$
      & $\Knows p \lif p$ \hfill \newline{(Vérité de la connaissance)} \\
      \hline
      \emph{transitive} : \newline
      $\forall u \forall v \forall w ((Ruv \land Rvw) \lif Ruw)$ &
      $\Knows p \lif \Knows \Knows p$ \hfill
           \newline{(Introspection positive)} \\
      \hline
      \emph{euclidienne} : \newline
      $\forall w \forall u \forall v ((Rwu \land Rwv) \lif Ruv)$ &
      $\lnot \Knows p \lif \Knows \neg \Knows p$ \hfill
        \newline{(Introspection négative)} \\
      \hline
    \end{tabular}
    \caption{Quatre principes épistémiques.}
    \ollabel{tab:four}
  \end{table}

Le principe de vérité de la connaissance, qui correspond à l'axiome
$T$, est souvent considéré comme le moins controversé de ces
principes : si !!a{formula} est connue, elle doit être vraie. La clôture ainsi que
l'introspection positive et l'introspection négative sont beaucoup
plus discutées.

La clôture, qui correspond à l'axiome $K$, exprime l'idée qu'un agent
connaissant une implication et son antécédent en connaît aussi le
conséquent. Cela peut paraître plausible dans certains cas. Par
exemple, je peux savoir que, si je suis à Victoria, je suis sur
l'île de Vancouver. Sauf à envisager d'étranges scénarios sceptiques,
je sais que je suis à Victoria ; je devrais donc savoir aussi que
je suis sur l'île de Vancouver. Dans ce cas, la clôture logique de
mes connaissances peut sembler assez intuitive. En revanche, nous
ne tirons pas toujours toutes les conséquences de ce que nous savons,
ce qui peut rendre le principe moins intuitif dans d'autres cas.

L'introspection positive, parfois appelée principe KK, affirme que,
si je sais quelque chose, je sais que je le sais. C'est le pendant
épistémique de l'axiome 4. L'introspection négative affirme que,
si je \emph{ne sais pas} quelque chose, je sais que je ne le sais pas ;
c'est le pendant de l'axiome 5. Ces deux principes semblent admettre des contre-exemples assez
ordinaires, dans lesquels je ne suis pas certain de savoir ou non
quelque chose que je sais en fait.
%
%

\end{document}
